\documentclass{article}
\usepackage{graphicx}
\usepackage{booktabs}
\usepackage[margin=2.5cm]{geometry}

\begin{document}

\title{ANUGA applied to flooding in Towradgi}

\maketitle

\section{Introduction}
This model simulates the Towradgi Creek 17 August 1998 flood. It was largely
developed by Petar Milevski in 2014 (with amendments in 2022), with some edits
by Gareth Davies and Steve Roberts. Other flood studies which model this event
can be found online from 2003 (Mike 11) and 2014 (TUFLOW) -- google `Towradgi
Creek Flood Study' for more information and to compare results across models.

The model is run with \texttt{anuga\_toml\_run towradgi.toml}, the
configuration generated by \texttt{make\_towradgi\_toml.py}. It uses about
257\,000 triangles, refined to 8\,m$^2$ along the creek banks, on a 1\,m DEM
with the bridge decks burned in. It includes 22 culverts and bridges as Boyd
box and pipe operators, handrail blockage as riverwalls, 49 gauge-polygon
rainfall series, a constant creek inflow of 3\,m$^3$/s at the western
boundary, and a fixed ocean level of 0\,m AHD. It does not include the piped
(minor) drainage system.

The run starts at $t = -3$\,h. During this burn-in only the creek inflow is
applied, which wets the channel before the rain starts at $t = 0$. All times
below are on the rainfall clock.

The figures and statistics below come from \texttt{Compare\_results\_with\_fieldObs.py},
run on the output of \texttt{towradgi.toml}.

\section{Results}

\subsection{Peak flood levels}

Figure~\ref{fig:spatial} shows the modelled peak flood depth and the error at
each of the 48 surveyed flood marks, and Figure~\ref{fig:hist} compares the
modelled and observed levels. Table~\ref{tab:stats} summarises the errors.

\begin{table}[h]
\centering
\begin{tabular}{lr}
\toprule
Flood marks & 48 \\
RMS error & 0.37\,m \\
Mean error (bias) & $-0.16$\,m \\
Mean absolute error & 0.31\,m \\
Largest error & 0.72\,m \\
Marks within 0.3\,m / 0.5\,m & 28 / 38 \\
\bottomrule
\end{tabular}
\caption{Modelled peak level minus surveyed flood level at the 48 marks.}
\label{tab:stats}
\end{table}

The model is slightly low overall. The largest under-predictions are in the
reach between the Railway Bridge and Pioneer Road (marks 1222 and 1224, about
$-0.7$\,m) and at a few marks on the side branches. Previous studies also
found this reach hard to match. The largest over-predictions are at the two
lagoon marks below Pioneer Road (1232 and 1234, $+0.36$ and $+0.54$\,m), where
the creek mouth in the DEM is flat at $-0.4$\,m and probably shallower than it
was during the flood. Lowering the mouth bed by 1\,m brings these two marks
down by about 0.15\,m without changing the others.

\begin{figure}
\centering
\includegraphics[width=\textwidth]{Spatial_Depth_and_Error.png}
\caption{Modelled peak depth, and the difference between the modelled peak
level and the surveyed flood level at each mark.}
\label{fig:spatial}
\end{figure}

\begin{figure}
\centering
\includegraphics[width=\textwidth]{Error_peakstage.png}
\caption{Modelled against observed flood levels, and the distribution of the
errors.}
\label{fig:hist}
\end{figure}

\subsection{Pioneer Road gauge}

Figure~\ref{fig:pioneer} shows the modelled and recorded stage at the Pioneer
Road bridge gauge. The logger clock runs 9\,h ahead of the rainfall clock, so
the record is plotted 9\,h earlier. The source of the offset is not
documented; it is the shift used since the original 2014 study, and fitting
the model to the record gives the same value. The model reproduces the timing
of the three main peaks. It reaches 4.10\,m against the recorded 4.32\,m at the
same time (11.25\,h), and it follows the falling limb closely (RMS error
0.13\,m from 13\,h to 24\,h; 0.31\,m over the whole storm).

The model rises later than the record in the first three hours of rain. Before
the storm the logger reads a flat 0.7\,m, while the modelled gauge cell is dry
until about 1.3\,h. That flat reading is probably the logger's dry floor or a
still pool rather than flowing water.

The 3\,m$^3$/s creek inflow replaces the 20\,m$^3$/s of the original script.
With 20\,m$^3$/s the falling limb stays about 0.7\,m above the record (gauge RMS
error 0.67\,m); the flood-mark errors are almost the same either way.

\begin{figure}
\centering
\includegraphics[width=\textwidth]{Pioneer_Bridge_Stage.png}
\caption{Modelled and recorded stage at the Pioneer Road bridge gauge. The
shaded band is the burn-in before the rain starts.}
\label{fig:pioneer}
\end{figure}

\subsection{Sensitivity}

Several variations each changed the flood-mark RMS error by less than
0.01\,m: a 25\% blockage of the bridges, raised deck walls on Pioneer Road, a
0.7\,m sill under the Pioneer Road bridge, a lagoon perched at 0.7\,m behind a
closed entrance, and the deeper creek mouth above. Removing all 22 structures
raises the RMS error to 0.45\,m and the largest error to 1.26\,m.

\end{document}
